% =============================================================================
% Chapter 2: Related Work (4-5 pages)
% =============================================================================

\section{Related Work}
\label{sec:related_work}

This chapter demonstrates the contribution of the thesis within the framework of four
areas of previous work: the formulation of community detection as a
quality-function optimization problem, QUBO-based approaches and
their hardware solvers, neural network methods for community
detection, and the unsupervised QUBO-via-GNN paradigm that the
thesis builds on. The chapter closes with a brief discussion of the
trivial collapse phenomenon, which motivates the regularization
analysis in later sections.

% -----------------------------------------------------------------------------
\subsection{Community detection problem}
\label{sec:rw_problem}

\subsubsection{Modularity and its limitations}
\label{sec:rw_modularity}

Community detection has been studied for over two decades; surveys
such as Fortunato~\cite{fortunato2010community} cover the main
threads. The dominant quality function is Newman's
modularity~\cite{newman2006modularity}, defined in
equation~\eqref{eq:modularity_intro}. Modularity compares the
observed edge density inside each community to the density expected
in a random graph with the same degree sequence, and high modularity
is interpreted as evidence of community structure.

Despite its widespread use, modularity has known issues. The
\emph{resolution limit} of Fortunato and
Barth\'elemy~\cite{fortunato2007resolution} shows that maximizing
modularity in large graphs systematically merges small communities
into larger ones, even when the small communities are well defined.
Modularity is also degenerate in the sense of Good et
al.~\cite{good2010performance}: many partitions of a real-world
graph achieve modularity close to the maximum, so the optimum is
neither unique nor robust. Alternative quality functions such as
normalized cut and conductance~\cite{shi2000normalized} avoid some
of these issues at the cost of different trade-offs. This thesis
keeps modularity as the optimization target, both because it is the
objective in the QUBO formulations of Wang et
al.~\cite{wang2024quantum} that the thesis adopts, and because
modularity remains the standard against which community detection
methods are typically compared.

\subsubsection{Classical algorithms}
\label{sec:rw_classical}

The Louvain algorithm of Blondel et al.~\cite{blondel2008louvain}
optimizes modularity through a two-phase greedy procedure. In the
first phase, vertices are repeatedly moved between communities to
greedily increase modularity; in the second phase, the resulting
communities are aggregated into super-vertices and the process is
repeated on the smaller graph. The Leiden algorithm of Traag et
al.~\cite{traag2019leiden} refines Louvain by adding a vertex-level
refinement step that guarantees connected communities and avoids
some pathological cases of Louvain. Both algorithms run in nearly
linear time on sparse graphs and are the de facto baselines for
modularity-based community detection.

Other classical approaches include spectral methods, which relax the
discrete partition problem to a continuous eigenvector problem on
the graph Laplacian or the modularity
matrix~\cite{newman2013spectral, shi2000normalized}, and
information-theoretic methods such as
Infomap~\cite{rosvall2008infomap}, which finds communities by
minimizing the description length of a random walk on the graph.
Stochastic block models~\cite{karrer2011stochastic} take a generative
view of the problem and recover communities through statistical
inference. The thesis uses Louvain as the primary baseline
because it shares the modularity objective with the proposed
method, making the comparison meaningful; Leiden is reported
alongside Louvain in the dataset-level baseline table
(Section~\ref{sec:exp_baselines_metrics}) as a sanity check and
behaves within a few thousandths of Louvain on every benchmark.

% -----------------------------------------------------------------------------
\subsection{QUBO formulations and hardware solvers}
\label{sec:rw_qubo}

\subsubsection{QUBO and Ising}
\label{sec:rw_lucas}

A Quadratic Unconstrained Binary Optimization (QUBO) problem
minimizes a quadratic form $\xb^T \Q \xb$ over binary vectors $\xb
\in \{0,1\}^N$. The QUBO problem is equivalent to the Ising spin
model under the change of variables $s_i = 2 x_i - 1$, mapping $\xb
\in \{0,1\}^N$ to $\mathbf{s} \in \{-1, +1\}^N$. This equivalence is
useful because the Ising model has a long history in physics and
because both quantum and analog hardware solvers are typically built
in the Ising form.

Lucas~\cite{lucas2014ising} provides the standard reference catalog
of NP-hard problems formulated as Ising/QUBO problems, including
maximum cut, graph coloring, traveling salesman, and others. This
catalog enables a single hardware solver to address a wide range of
combinatorial problems by translating each problem into the
hardware's native form.

\subsubsection{QUBO for community detection}
\label{sec:rw_qubo_cd}

The use of QUBO for community detection traces back to the
graph-partitioning formulation of Ushijima-Mwesigwa et
al.~\cite{ushijima2017graph}, who solved a balanced 2-way partition
problem on the D-Wave quantum annealer. Negre et
al.~\cite{negre2020detecting} extended the approach to multi-way
community detection by introducing one-hot encoded community
indicators with a soft constraint penalty, which is the
multi-class formulation later adopted by Wang et
al.~\cite{wang2024quantum}.

The most direct precursor of this thesis is the work of Wang et
al.~\cite{wang2024quantum}, which formulates modularity maximization
as both a bipartition QUBO (formula~6 in their paper, used here as
the basis of Section~\ref{sec:method_qubo_k2}) and a multi-class
QUBO (formula~5, used in Section~\ref{sec:method_qubo_multi}). They
solve the resulting QUBO instances on a Coherent Ising Machine
(CIM) and apply the method to a real banking fraud detection
dataset. Their results demonstrate that QUBO-based community
detection can match or exceed classical heuristics in some settings,
but the underlying solver depends entirely on specialized analog
hardware.

\subsubsection{Hardware solvers and their limitations}
\label{sec:rw_hardware}

Hardware QUBO solvers fall into two main families. Quantum annealers
such as the D-Wave systems implement the Ising model on a network
of physical qubits and find low-energy states through quantum
tunneling. Coherent Ising Machines, used by Wang et
al.~\cite{wang2024quantum}, implement the Ising model through
optical parametric oscillators and rely on synchronization dynamics
rather than quantum tunneling. Both families share two practical
limitations. First, the available system size is bounded by the
hardware's number of qubits or oscillators, with current systems
supporting at most a few thousand variables. For multi-class
community detection this corresponds to graphs with at most a few
hundred vertices when $k$ is moderate. Second, both families require
specialized hardware that is not widely accessible.

These limitations motivate the search for software-based QUBO
solvers that run on commodity hardware while preserving the global
nature of the QUBO formulation. The unsupervised neural approach
discussed in Section~\ref{sec:rw_unsupervised} is one such direction.

% -----------------------------------------------------------------------------
\subsection{GNN-based community detection}
\label{sec:rw_gnn_cd}

\subsubsection{Graph neural networks}
\label{sec:rw_gnn_arch}

Graph neural networks (GNNs) operate on graph-structured data by
iteratively updating vertex embeddings based on the embeddings of
neighboring vertices. The graph convolutional network (GCN) of Kipf
and Welling~\cite{kipf2017semi} popularized this class of models;
GraphSAGE~\cite{hamilton2017inductive} introduced sampled and
concatenation-based aggregation; the graph attention network
(GAT)~\cite{velickovic2018graph} added learned attention weights
between neighbors; and GIN~\cite{xu2019gin} provided theoretical
guarantees on representational power. The methodology in this
thesis builds on a GraphSAGE backbone, following the choice made by
the QIGNN framework discussed in Section~\ref{sec:rw_qignn}.

\subsubsection{Supervised approaches: DMoN}
\label{sec:rw_dmon}

Tsitsulin et al.~\cite{tsitsulin2020graph} proposed the
Differentiable Modularity Optimization Network (DMoN), which uses
a GNN to map vertex features to soft community assignments and
trains the GNN end-to-end by maximizing a differentiable surrogate
of modularity. DMoN is supervised in the sense that it requires
node features (typically derived from labeled data) but is
unsupervised in its training signal: only the modularity-based loss
is used, with no community labels.

DMoN includes an orthogonality regularization that penalizes
overlap between community assignments. The specific orthogonality
formula
$\bigl\| \Pb^T \Pb / \| \Pb^T \Pb \|_F - I_k / \sqrt{k} \bigr\|_F^2$
adopted in Section~\ref{sec:method_ortho} of this thesis is the
form introduced by MinCutPool~\cite{bianchi2020spectral} and also
used by DMoN; both papers report it as effective against
degenerate cluster assignments. This thesis applies it in the
rather different setting of an unsupervised QUBO loss.

A related supervised method is MinCutPool by Bianchi et
al.~\cite{bianchi2020spectral}, which uses a normalized minimum
cut as the training signal and includes a similar orthogonality
regularization. Both methods rely on supervised features and are
designed for graph-level pooling rather than full community
detection on a single graph.

% -----------------------------------------------------------------------------
\subsection{Unsupervised neural QUBO solvers}
\label{sec:rw_unsupervised}

\subsubsection{PI-GNN}
\label{sec:rw_pignn}

The Physics-Inspired GNN (PI-GNN) of Schuetz et
al.~\cite{schuetz2022combinatorial} is the starting point for the
line of work this thesis extends. Given a QUBO instance defined by
a matrix $\Q$, PI-GNN trains a GNN whose output is a vector of
relaxed assignments $\pb \in [0,1]^N$, and the GNN is trained on a
single graph by minimizing the relaxed QUBO loss $\pb^T \Q \pb$,
without any labeled data. After training, the relaxed solution is
rounded to a binary vector. The approach is validated on Maximum
Independent Set and MaxCut benchmarks, where PI-GNN matches or
exceeds classical heuristics on a number of instances while running
on a single GPU.

The key insight of Schuetz et al.\ is that the GNN provides an
implicit regularization through the graph structure: vertices that
are close in the graph receive similar embeddings, which translates
to similar relaxed assignments. This makes the optimization easier
than direct gradient descent on $\pb$ in $\mathbb{R}^N$, where no
such structure is exploited.

\subsubsection{QIGNN}
\label{sec:rw_qignn}

The QIGNN framework~\cite{qignn2026} extends PI-GNN with iterative
refinement. At each training epoch, the network's previous-epoch
output is concatenated with its static input features, so that the
GNN can revisit and improve its earlier prediction. The
architecture itself uses GraphSAGE blocks with two parallel
aggregations (mean and max-pool) combined through a residual
connection, giving a richer neighborhood representation than a
single-aggregator GraphSAGE block. The framework is validated on
the same MaxCut and MIS benchmarks as PI-GNN and reports improved
stability and final solution quality.

QIGNN is the architectural starting point of this thesis. The
existing validation of QIGNN is restricted to two-class
combinatorial problems, where each variable is a single binary
indicator; the present thesis is, to the best of the author's
knowledge, the first work to apply the QIGNN paradigm to a
multi-class problem.

\subsubsection{Limits of neural QUBO solvers}
\label{sec:rw_critique}

The unsupervised neural QUBO solver paradigm has been challenged.
Angelini and Ricci-Tersenghi~\cite{angelini2023modern} argue that
on Maximum Independent Set the original PI-GNN is outperformed by a
simple greedy degree-based heuristic, and that the apparent
strength of PI-GNN in the original paper is partly due to the choice
of benchmark instances. Their critique does not invalidate the
general approach, but it does highlight the importance of careful
empirical evaluation against strong classical baselines, rather than
adoption of neural QUBO solvers based on a small set of benchmark
results. The empirical evaluation in this thesis follows this
guidance by including Louvain as the primary baseline (with Leiden as a sanity-check secondary baseline) on a
diverse collection of graphs.

% -----------------------------------------------------------------------------
\subsection{Trivial collapse and regularization}
\label{sec:rw_collapse}

The trivial collapse phenomenon~--~when a clustering model assigns
all vertices to a single cluster~--~is a recurring difficulty in
unsupervised methods. Tsitsulin et al.~\cite{tsitsulin2020graph}
and Bianchi et al.~\cite{bianchi2020spectral} report it in GNN-based
clustering and address it through orthogonality and collapse
regularizations.

These regularizations have been studied in supervised settings with
node features and a direct modularity surrogate. The QUBO-based
unsupervised setting considered here differs in both respects: there
are no input features beyond the graph structure, and the loss is a
QUBO form with a soft constraint penalty rather than a direct
modularity surrogate. The behavior of these regularizations in this
setting has not, to the author's knowledge, been previously studied;
Section~\ref{sec:exp_regularization} provides an empirical analysis
across graphs of varying balance and number of communities.

\newpage